\section{Dirichlet's and Neumann's problems on a plane}
\label{sec:dirichlet_neumann}

\subsection{Horizontal plane $z^{\dagger}_{0}$ located above the sources}
\label{subsec:dirichlet_neumann_above_sources}

Let $V(\mathbf{r})$ and $W(\mathbf{r})$ be surface integrals given by:
\begin{equation}
V(\mathbf{r}) =  
\int\limits_{-\infty}^{\infty} \int\limits_{-\infty}^{\infty}
U(\mathbf{r}^{\dagger}) \:\: 
\partial_{z} \Psi(\mathbf{r}, \mathbf{r}^{\dagger}) 
\: dx^{\dagger}dy^{\dagger} \: , \quad z < z^{\dagger}_{0} < z^{t} \: ,
\label{eq:scalar-potential-V}
\end{equation}
and 
\begin{equation}
W(\mathbf{r}) =  
\int\limits_{-\infty}^{\infty} \int\limits_{-\infty}^{\infty}
\partial_{z} U(\mathbf{r}^{\dagger}) \:\: 
\Psi(\mathbf{r}, \mathbf{r}^{\dagger}) 
\: dx^{\dagger}dy^{\dagger} \: , \quad z < z^{\dagger}_{0} < z^{t} \: ,
\label{eq:scalar-potential-W}
\end{equation}
where $z^{t}$ defines the top of the sources, 
$\mathbf{r} = (x, y, z)^{\top} \notin \overline{\mathcal{V}}$ is an observation point,
$\partial_{z} U(\mathbf{r}^{\dagger}) \equiv \tfrac{\partial U(\mathbf{r}^{\dagger})}{\partial z}$
is the vertical derivative of the potential $U(\mathbf{r}^{\dagger})$
(equation \ref{eq:1st-deriv-U}) defined at a point 
$\mathbf{r}^{\dagger} = (x^{\dagger}, y^{\dagger}, z^{\dagger}_{0})^{\top} \notin \overline{\mathcal{V}}$ 
on the horizontal plane $z^{\dagger}_{0} < z^{t}$ located above the sources, and 
$\partial_{z} \Psi(\mathbf{r}, \mathbf{r}^{\dagger}) \equiv \tfrac{\partial
\Psi(\mathbf{r}, \mathbf{r}^{\dagger})}{\partial z}$ is the vertical derivative of 
\begin{equation}
\Psi(\mathbf{r}, \mathbf{r}^{\dagger}) = 
\frac{1}{2\pi} \: 
\frac{1}{\| \mathbf{r} - \mathbf{r}^{\dagger} \|} \: ,
\label{eq:Psi}
\end{equation}
which represents the scaled inverse distance function between 
$\mathbf{r}$ and $\mathbf{r}^{\dagger}$.

It can be shown that, at all points $\mathbf{r}$ located outside the sources (i.e., 
$\mathbf{r} \notin \overline{\mathcal{V}}$), $U(\mathbf{r})$ (equation
\ref{eq:scalar-potential-U}) has continuous partial derivatives of all orders
and satisfies the Laplace's equation \citep[][ p. 122--124]{kellogg1967}.
Besides, it is regular at infinity \citep[][ p. 217]{kellogg1967}.
Because of these properties, we have that
\begin{equation}
	U(\mathbf{r}) = V(\mathbf{r}) \: , \quad z < z^{\dagger}_{0} < z^{t} \: ,
	\label{eq:U-Dirichlet}
\end{equation}
and
\begin{equation}
	U(\mathbf{r}) = W(\mathbf{r}) \: , \quad z < z^{\dagger}_{0} < z^{t} \: .
	\label{eq:U-Neumann}
\end{equation}
The common outline to deduce equations \ref{eq:U-Dirichlet} and \ref{eq:U-Neumann}
can be found, for example, in \citet{roy1962}.

Equation \ref{eq:U-Dirichlet} states that $U(\mathbf{r})$ can be defined
in terms of its values $U(\mathbf{r}^{\dagger})$ on the plane $z^{\dagger}_{0}$ and 
represents the solution of the \textit{Dirichlet’s problem} or the \textit{first boundary
value problem of potential theory} \cite[][ p. 236]{kellogg1967} on a plane.
On the other hand, equation \ref{eq:U-Neumann} shows that $U(\mathbf{r})$ can 
be defined in terms of its vertical derivative $\partial_{z} U(\mathbf{r}^{\dagger})$ 
on the plane $z^{\dagger}_{0}$ and represents the solution of the 
\textit{Neumann’s problem} or the \textit{second boundary value problem of potential
theory} \cite[][ p. 246]{kellogg1967} on a plane.

It is important to stress that equations \ref{eq:U-Dirichlet} and \ref{eq:U-Neumann}
are valid for the case in which (i) $\mathbf{r} = (x, y, z)^{\top}$ represents any observation 
point located above the sources, not necessarily on a horizontal plane, and (ii)  
$\mathbf{r}^{\dagger} = (x^{\dagger}, y^{\dagger}, z^{\dagger}_{0})^{\top}$ represents 
points on the horizontal plane $z^{\dagger}_{0}$ located above the sources, 
so that $z < z^{\dagger}_{0} < z^{t}$. What follows aims at generalizing
these equations for the case in which they are defined on a horizontal plane that
crosses the sources. 

\subsection{Horizontal plane $z^{\ddagger}_{0}$ that crosses the sources}
\label{subsec:generalized_dirichlet_neumann}

\subsubsection{Approximated Dirichlet's problem}

Let $V_{\varepsilon}(\mathbf{r}^{\dagger})$ be the surface integral given by:
\begin{equation}
V_{\varepsilon}(\mathbf{r}^{\dagger}) = 
\int\limits_{-\infty}^{\infty} \int\limits_{-\infty}^{\infty}
\tau_{\varepsilon}(\mathbf{r}^{\ddagger}) \:
\partial_{z} \Psi(\mathbf{r}^{\dagger}, \mathbf{r}^{\ddagger}) 
\:\: dx^{\ddagger}dy^{\ddagger} \: , \quad z^{\dagger}_{0} < z^{t} < z^{\ddagger}_{0} \: ,
\label{eq:V-tau-epsilon}
\end{equation}
where 
$\mathbf{r}^{\ddagger} = (x^{\ddagger}, y^{\ddagger}, z^{\ddagger}_{0})^{\top}$ is
a point on the horizontal plane that crosses the 3D sources at $z^{\ddagger}_{0} > z^{t}$. 
Our goal here is finding an appropriate function $\tau_{\varepsilon}$ 
(equation \ref{eq:V-tau-epsilon}) so that
\begin{equation}
	\int\limits_{-\infty}^{\infty} \int\limits_{-\infty}^{\infty}
	\vert U(\mathbf{r}^{\dagger}) - 
	V_{\varepsilon}(\mathbf{r}^{\dagger}) \vert^{2}
	\:\: dx^{\dagger}dy^{\dagger} \le \varepsilon \: ,
	\label{eq:epsilon-bound}
\end{equation}
where $\varepsilon > 0$ is a fixed but arbitrarily small error level.
This can be made by first applying the Fourier-convolution theorem 
\cite[][ p. 268]{blakely1996} to $V_{\varepsilon}(\mathbf{r}^{\dagger})$ 
(equation \ref{eq:V-tau-epsilon}) to obtain:
\begin{equation}
	\mathcal{F} [ V_{\varepsilon}(\mathbf{r}^{\dagger}) ]
	= \mathcal{F} [ \tau_{\varepsilon}(\mathbf{r}^{\ddagger}) ] \:
	\mathcal{F} [ \partial_{z} \Psi(\mathbf{r}^{\dagger}, \mathbf{r}^{\ddagger}) ] \: ,
	\label{eq:Fourier-V-tau-epsilon}
\end{equation}
where
\begin{equation}
	\mathcal{F} [ \partial_{z} \Psi(\mathbf{r}^{\dagger}, \mathbf{r}^{\ddagger}) ] 
	= e^{- k \, (z^{\ddagger}_{0} - z^{\dagger}_{0})}
	\: , \quad z^{\ddagger}_{0} > z^{\dagger}_{0} \: , \quad k \ne 0 \: ,
	\label{eq:Fourier-Dz-Psi}
\end{equation}
is defined according to \citet[][ p. 275]{blakely1996}, with 
$k = \sqrt{k_{x}^{2} + k_{y}^{2}}$. 
Now, we define a general solution for $\tau_{\varepsilon}$ as follows:
\begin{equation}
	\mathcal{F} \left[ \tau_{\varepsilon}(\mathbf{r}^{\ddagger}) \right] = 
	\begin{cases}
		\mathcal{F} \left[ U(\mathbf{r}^{\dagger}) \right] 
		e^{\, k \, (z^{\ddagger}_{0} - z^{\dagger}_{0})} & , 
		\quad 0 < k \le k_{\varepsilon} \: , \\
		0 & , \quad k > k_{\varepsilon} \: ,
	\end{cases}
	\label{eq:Fourier-tau-epsilon}
\end{equation}
where $k_{\varepsilon}$ is a fixed and possibly very large truncating wavenumber.

Note that $\tau_{\varepsilon}(\mathbf{r}^{\ddagger})$ 
approximates the potential field obtained by computing the downward continuation of 
$U(\mathbf{r}^{\dagger})$ (equation \ref{eq:scalar-potential-U}) 
to points located inside the sources. 
It is well-known, however, that downward continuation can only be computed at 
source-free points, where the Laplace's equation applies. This means that the
continuation plane is limited by the top $z^{t}$ of the sources, otherwise the continued field 
presents an oscillatory character \citep[e.g.,][]{roy1967, fedi_florio2002} that may
prevent $\mathcal{F} \left[ \tau_{\varepsilon}(\mathbf{r}^{\ddagger}) \right]$ 
to be square integrable.
In this case, the Fourier-convolution theorem cannot be applied.
The truncating wavenumber $k_{\varepsilon}$ (equation \ref{eq:Fourier-tau-epsilon}) 
overcomes this problem and guarantees that the Fourier-convolution can be applied to
$V_{\varepsilon}(\mathbf{r}^{\dagger})$ (equation \ref{eq:V-tau-epsilon}). 

We get by substituting 
$\mathcal{F} \left[ \partial_{z} \Psi(\mathbf{r}^{\dagger}, \mathbf{r}^{\ddagger}) \right]$
(equation \ref{eq:Fourier-Dz-Psi}) and 
$\mathcal{F} \left[ \tau_{\varepsilon}(\mathbf{r}^{\ddagger}) \right]$
(equation \ref{eq:Fourier-tau-epsilon}) into equation \ref{eq:Fourier-V-tau-epsilon} that
\begin{equation}
	\mathcal{F} [ V_{\varepsilon}(\mathbf{r}^{\dagger}) ] = 
	\begin{cases}
		\mathcal{F} [ U(\mathbf{r}^{\dagger}) ] &
		, \quad 0 < k \le k_{\varepsilon} \: , \\
		0 & , \quad k > k_{\varepsilon} \: ,
	\end{cases}
	\label{eq:Fourier-V-tau-epsilon-2}
\end{equation}
which shows that the Fourier transforms of 
$U(\mathbf{r}^{\dagger})$ (equation \ref{eq:scalar-potential-U}) and 
$V_{\varepsilon}(\mathbf{r}^{\dagger})$ (equation \ref{eq:V-tau-epsilon}) 
on the horizontal plane $z^{\dagger}_{0}$ located above the 3D sources coincide up 
to the truncating wavenumber $k_{\varepsilon}$.
Then, by applying the Plancherel's theorem 
\citep[][ p. 119]{bracewell1999} to equation \ref{eq:epsilon-bound} and
changing variables with $k_{x} = k \, \cos\theta$, 
$k_{y} = k \, \sin\theta$ and $k = \sqrt{k_{x}^{2} + k_{y}^{2}}$ we get
\begin{equation}
	\int\limits_{-\infty}^{\infty} \int\limits_{-\infty}^{\infty}
	\vert 
	U(\mathbf{r}^{\dagger}) - V_{\varepsilon}(\mathbf{r}^{\dagger})
	\vert^{2} 
	\: dx^{\dagger}dy^{\dagger} = 
	\int\limits_{0}^{2\pi} \int\limits_{0}^{\infty}
	\vert \mathcal{F} [ U(\mathbf{r}^{\dagger}) ] - \mathcal{F} [
	V_{\varepsilon}(\mathbf{r}^{\dagger}) ] \vert^{2} \: k \, dkd\theta 	
	\quad .
	\label{eq:Plancherel-epsilon-bound}
\end{equation}
If we now substitute equation \ref{eq:Fourier-V-tau-epsilon-2} into
\ref{eq:Plancherel-epsilon-bound}, we obtain
\begin{equation}
	\int\limits_{-\infty}^{\infty} \int\limits_{-\infty}^{\infty}
	\vert 
	U(\mathbf{r}^{\dagger}) - V_{\varepsilon}(\mathbf{r}^{\dagger})
	\vert^{2} 
	\: dx^{\dagger}dy^{\dagger} = 
	\int\limits_{0}^{2\pi} \int\limits_{k_{\varepsilon}}^{\infty}
	\vert \mathcal{F} [ U(\mathbf{r}^{\dagger}) ] \vert^{2} \: k \, dkd\theta 	
	\quad .
	\label{eq:Plancherel-epsilon-bound2}
\end{equation}

Because we know that $\mathcal{F} \left[ U(\mathbf{r}^{\dagger}) \right]$ 
(equation \ref{eq:Fourier-U-r-dagger}) is square integrable 
(equation \ref{eq:abs-Fourier-U-r-dagger-bound}), we may use an adapted version 
of the Runge–Walsh theorem 
\citep[e.g.,][ p. 152]{freeden_gerhards2012} to affirm that there exists a sufficiently
large $k_{\varepsilon}$ such that the right hand side of 
\ref{eq:Plancherel-epsilon-bound2} is smaller than or equal to a predefined 
$\varepsilon > 0$, which guarantees that \ref{eq:epsilon-bound} holds true.

PAREI AQUI

It is important to stress, however, that this is a pure existence theorem and 
we are generally unable to provide a particular rule on how to choose $k_{\varepsilon}$.
Fortunately, we may substitute the upper bound for the amplitude spectrum of
$\mathcal{F} [ U(\mathbf{r}^{\dagger}) ]$ (equation \ref{eq:abs-Fourier-U-r-dagger-bound})
into the right side of equation \ref{eq:Plancherel-U-tau-epsilon-2} to define an 
$\varepsilon$ (equation \ref{eq:epsilon}) as follows:
\begin{equation}
	\int\limits_{0}^{2\pi} \int\limits_{k_{\varepsilon}}^{\infty}
	\vert \mathcal{F} [ U(\mathbf{r}^{\dagger}) ] \vert^{2} \: k \, dkd\theta 
	\le 2\pi \, \Omega^{2} \, \Gamma(0, 2 \, \Delta z \, k_{\varepsilon}) = \varepsilon \quad ,
	\label{eq:epsilon-Gamma-function}
\end{equation}
where 
\begin{equation}
	\Gamma(0, 2 \, \Delta z \, k_{\varepsilon}) = \int\limits_{k_{\varepsilon}}^{\infty}
	\frac{e^{- 2 \, \Delta z \, k }}{k} \: dk \: , \quad \Delta z = z^{t} - z^{\dagger}_{0} > 0 \: ,
	\label{eq:Gamma-function}
\end{equation}
is the Incomplete Gamma function \citep[][ p. 899]{gradshteyn_ryzhik2007}.

We have just shown that there exists a function $\tau_{\varepsilon}(\mathbf{r}'')$
(equation \ref{eq:Fourier-tau-epsilon}) defining the potential 
$U_{\varepsilon}(\mathbf{r}^{\dagger})$ (equation \ref{eq:U-tau-epsilon}), which
reproduces $U(\mathbf{r}^{\dagger})$ (equation \ref{eq:scalar-potential-U}) on the
horizontal plane $z^{\dagger}_{0}$ above the sources to some predefined error 
level $\varepsilon$ (equations \ref{eq:epsilon} and \ref{eq:epsilon-Gamma-function}).
This result allows us generalizing the Dirichlet's problem 
(equation \ref{eq:U-Dirichlet}) as follows:
The potential field $U(\mathbf{r})$ (equation \ref{eq:scalar-potential-U}) is exactly reproduced
in terms of its values on a plane $z^{\dagger}_{0}$ located above the sources,
according to equation \ref{eq:U-Dirichlet}, or approximated by $U_{\varepsilon}$ 
(equation \ref{eq:U-tau-epsilon}) to a predefined error level $\varepsilon$
(equations \ref{eq:epsilon} and \ref{eq:epsilon-Gamma-function}) in terms of
a continuous function $\tau_{\varepsilon}$ (equation \ref{eq:Fourier-tau-epsilon}).


%What follows aims at investigating if this is also the case for points located 
%above $z^{\dagger}_{0}$ (i.e., at points with vertical coordinate $z < z^{\dagger}_{0}$).
%To do this, let us first define the solution of the Dirichlet's problem 
%(equation \ref{eq:U-Dirichlet}) for the harmonic function 
%$U_{\varepsilon}(\mathbf{r}^{\dagger})$ (equation \ref{eq:U-tau-epsilon}) on the
%plane $z^{\dagger}_{0}$ and change the integration order:
%\begin{equation}
%\begin{split}
%U_{\varepsilon}(\mathbf{r}) 
%&= \int\limits_{-\infty}^{\infty} \int\limits_{-\infty}^{\infty}
%\left[ 
%\int\limits_{-\infty}^{\infty} \int\limits_{-\infty}^{\infty}
%\tau_{\varepsilon}(\mathbf{r}'') \:\: 
%\partial_{z} \Psi(\mathbf{r}^{\dagger}, \mathbf{r}'') 
%\: dx''dy'' \right]
%\partial_{z} \Psi(\mathbf{r}, \mathbf{r}^{\dagger}) 
%\:\: dx^{\dagger}dy^{\dagger} \\
%&= 
%\int\limits_{-\infty}^{\infty} \int\limits_{-\infty}^{\infty}
%\tau_{\varepsilon}(\mathbf{r}'') \:\:
%\left[ 
%\int\limits_{-\infty}^{\infty} \int\limits_{-\infty}^{\infty}
%\partial_{z} \Psi(\mathbf{r}^{\dagger}, \mathbf{r}'')
%\: \partial_{z} \Psi(\mathbf{r}, \mathbf{r}^{\dagger}) 
%\: dx^{\dagger}dy^{\dagger} \right]
%\:\: dx''dy'' \quad .
%\end{split}
%\label{eq:intermediary-equation}
%\end{equation}
%Then, by noting that the term within brackets in the second line of equation
%\ref{eq:intermediary-equation} represents the solution of Dirichlet's problem
%for the function $\partial_{z} \Psi(\mathbf{r}, \mathbf{r}'')$ on the plane
%$z^{\dagger}_{0}$, we obtain:
%\begin{equation}
%U_{\varepsilon}(\mathbf{r}) =  
%\int\limits_{-\infty}^{\infty} \int\limits_{-\infty}^{\infty}
%\tau_{\varepsilon}(\mathbf{r}'') \:\: 
%\partial_{z} \Psi(\mathbf{r}, \mathbf{r}'') 
%\:\: dx''dy'' \: , \quad z < z''_{0} \: .
%\label{eq:U-tau-epsilon-Dirichlet-cross}
%\end{equation}
%Now, consider that both $U(\mathbf{r})$ (equation \ref{eq:U-Dirichlet}) and 
%$U_{\varepsilon}(\mathbf{r})$ (equation \ref{eq:U-tau-epsilon-Dirichlet-cross}) 
%are evaluated at the same points 
%$\mathbf{r}^{\ddagger} = (x^{\ddagger}, y^{\ddagger}, z^{\ddagger}_{0})^{\top}
%\notin \overline{\mathcal{V}}$, 
%on a horizontal plane located at $z^{\ddagger}_{0} < z^{\dagger}_{0}$.
%In this case, the difference 
%
%PAREI AQUI
%
%we integrate the squared absolute difference between $U(\mathbf{r})$ 
%(equation \ref{eq:U-Dirichlet}) and $U_{\varepsilon}(\mathbf{r})$ 
%(equation \ref{eq:U-tau-epsilon-Dirichlet-cross}) on a horizontal plane 
%located at $z^{\ddagger}_{0} < z^{\dagger}_{0}$ to obtain
%\begin{equation}
%	\int\limits_{-\infty}^{\infty} \int\limits_{-\infty}^{\infty}
%	\vert U(\mathbf{r}^{\ddagger}) - U_{\varepsilon}(\mathbf{r}^{\ddagger}) \vert^{2} 
%	\: dx^{\ddagger}dy^{\ddagger} 
%	= \int\limits_{-\infty}^{\infty} \int\limits_{-\infty}^{\infty}
%	\vert U(\mathbf{r}^{\dagger}) - U_{\varepsilon}(\mathbf{r}^{\dagger}) \vert^{2} \: 
%	\partial_{z} \Psi(\mathbf{r}^{\ddagger}, \mathbf{r}^{\dagger})^{2}
%	\: dx^{\dagger}dy^{\dagger} \: , \quad z^{\ddagger}_{0} < z^{\dagger}_{0} \: .
%	\label{eq:U-tau-epsilon-difference}
%\end{equation}
%Finally, we use the inequality 
%\begin{equation}
%	z^{\ddagger}_{0} \le z^{\dagger}_{0} - \sqrt{\tfrac{1}{2\pi}} \implies
%	\partial_{z} \Psi(\mathbf{r}^{\ddagger}, \mathbf{r}^{\dagger}) \le 1 \: ,
%	\label{eq:Dz-Phi-constraint}
%\end{equation}
%and equation \ref{eq:epsilon} to show that the left side of equation
%\ref{eq:U-tau-epsilon-difference} satisfies:
%\begin{equation}
%	\int\limits_{-\infty}^{\infty} \int\limits_{-\infty}^{\infty}
%	\vert U(\mathbf{r}^{\ddagger}) - U_{\varepsilon}(\mathbf{r}^{\ddagger}) \vert^{2} 
%	\: dx^{\ddagger}dy^{\ddagger} \le \varepsilon \: ,
%	\quad z^{\ddagger}_{0} \le z^{\dagger}_{0} - \sqrt{\tfrac{1}{2\pi}} \: .
%	\label{eq:U-error-level-r}
%\end{equation}

%PAREI AQUI 1
%
%Consider now the following integral equation:
%\begin{equation}
%	U_{\zeta}(\mathbf{r}^{\dagger}) = 
%	\int\limits_{-\infty}^{\infty} \int\limits_{-\infty}^{\infty}
%	\upsilon_{\zeta}(\mathbf{r}'') \:
%	\Psi(\mathbf{r}^{\dagger}, \mathbf{r}'') 
%	\:\: dx''dy'' \: ,
%	\label{eq:U-upsilon-zeta}
%\end{equation}
%subject to
%\begin{equation}
%	\int\limits_{-\infty}^{\infty} \int\limits_{-\infty}^{\infty}
%	\vert U(\mathbf{r}^{\dagger}) - 
%	U_{\zeta}(\mathbf{r}^{\dagger}) \vert^{2}
%	\:\: dx^{\dagger}dy^{\dagger} \le \zeta \: ,
%	\label{eq:zeta}
%\end{equation}
%where $\zeta > 0$ is a fixed but arbitrarily small error level, 
%$\Psi(\mathbf{r}^{\dagger}, \mathbf{r}'')$ is the scaled inverse distance function 
%(equation \ref{eq:Psi}) evaluated at the points
%$\mathbf{r}^{\dagger} = (x^{\dagger}, y^{\dagger}, z^{\dagger}_{0})^{\top}$ on the
%horizontal plane $z^{\dagger}_{0}$ above the 3D sources and 
%$\mathbf{r}'' = (x'', y'', z''_{0})^{\top}$ 
%on a horizontal plane that crosses the 3D sources at $z''_{0} = z^{\dagger}_{0} + \Delta z$, 
%with $\Delta z > 0$. 
%
%The same approach followed to find a function $\tau_{\varepsilon}$ 
%satisfying equation \ref{eq:U-tau-epsilon} can be applied to find a function
%$\tau_{\zeta}$ satisfying equation \ref{eq:U-upsilon-zeta}.
%We first apply the Fourier-convolution theorem 
%\cite[][ p. 268]{blakely1996} to $U_{\zeta}(\mathbf{r}^{\dagger})$ 
%(equation \ref{eq:U-upsilon-zeta}) and obtain:
%\begin{equation}
%	\mathcal{F} [ U_{\zeta}(\mathbf{r}^{\dagger}) ]
%	= \mathcal{F} [ \upsilon_{\zeta}(\mathbf{r}'') ] \:
%	\mathcal{F} [ \Psi(\mathbf{r}^{\dagger}, \mathbf{r}'') ] \: ,
%	\label{eq:Fourier-U-upsilon}
%\end{equation}
%where
%\begin{equation}
%	\mathcal{F} [ \Psi(\mathbf{r}^{\dagger}, \mathbf{r}'') ] 
%	= \frac{e^{- k \, (z''_{0} - z^{\dagger}_{0})}}{k}
%	\: , \quad z''_{0} > z^{\dagger}_{0} \: , \quad k \ne 0 \: ,
%	\label{eq:Fourier-Psi}
%\end{equation}
%is defined according to \citet[][ p. 275]{blakely1996}, with 
%$k = \sqrt{k_{x}^{2} + k_{y}^{2}}$. 
%Then we define a general solution for $\upsilon_{\zeta}$ as follows:
%\begin{equation}
%	\mathcal{F} \left[ \upsilon_{\zeta}(\mathbf{r}'') \right] = 
%	\begin{cases}
%		\mathcal{F} \left[ \partial_{z} U(\mathbf{r}^{\dagger}) \right] 
%		e^{\, k \, (z''_{0} - z^{\dagger}_{0})} & , 
%		\quad 0 < k \le k_{\zeta} \: , \\
%		0 & , \quad k > k_{\zeta} \: ,
%	\end{cases}
%	\label{eq:Fourier-upsilon-zeta}
%\end{equation}
%where $\mathcal{F} \left[ \partial_{z} U(\mathbf{r}^{\dagger}) \right] = 
%k \, \mathcal{F} \left[ U(\mathbf{r}^{\dagger}) \right]$
%\cite[][ p. 326]{blakely1996} and $k_{\zeta}$ is a fixed and possibly very 
%large truncating wavenumber.
%In this case, $\upsilon_{\zeta}(\mathbf{r}'')$ 
%approximates the potential field obtained by computing the downward continuation of 
%$\partial_{z} U(\mathbf{r}^{\dagger})$ (equation \ref{eq:U-Neumann}) 
%to points located inside the sources. 
%
%Then, by substituting 
%$\mathcal{F} \left[ \Psi(\mathbf{r}^{\dagger}, \mathbf{r}'') \right]$
%(equation \ref{eq:Fourier-Psi}) and 
%$\mathcal{F} \left[ \upsilon_{\zeta}(\mathbf{r}'') \right]$
%(equation \ref{eq:Fourier-upsilon-zeta}) into equation \ref{eq:Fourier-U-upsilon},
%as well using the identity 
%$\mathcal{F} \left[ \partial_{z} U_{\zeta}(\mathbf{r}^{\dagger}) \right] = 
%k \, \mathcal{F} \left[ U_{\zeta}(\mathbf{r}^{\dagger}) \right]$
%\cite[][ p. 326]{blakely1996}, we show that
%\begin{equation}
%	\mathcal{F} [ \partial_{z} U_{\zeta}(\mathbf{r}^{\dagger}) ] = 
%	\begin{cases}
%		\mathcal{F} [ \partial_{z} U(\mathbf{r}^{\dagger}) ] &
%		, \quad 0 < k \le k_{\zeta} \: , \\
%		0 & , \quad k > k_{\zeta} \: ,
%	\end{cases}
%	\label{eq:Fourier-Dz-U-upsilon-zeta}
%\end{equation}
%which shows that the Fourier transforms of 
%$\partial_{z} U(\mathbf{r}^{\dagger})$ (equation \ref{eq:1st-deriv-U}) and  
%the vertical derivative of $U_{\zeta}(\mathbf{r}^{\dagger})$ 
%(equation \ref{eq:U-upsilon-zeta}) 
%on the horizontal plane $z^{\dagger}_{0}$ located above the 3D sources coincide up 
%to the truncating wavenumber $k_{\zeta}$.
%
%
%Again, we may use the Plancherel's theorem \citep[][ p. 119]{bracewell1999} and 
%an adapted version of the Runge–Walsh theorem 
%\citep[e.g.,][ p. 152]{freeden_gerhards2012} to affirm that there exists a sufficiently
%large $k_{\zeta}$ such that equations \ref{eq:U-upsilon-zeta}
%and \ref{eq:zeta} hold true.
%
%Equations \ref{eq:U-upsilon-zeta}, \ref{eq:zeta} and \ref{eq:Fourier-upsilon-zeta}
%generalize the solution of Neumann's problem (equation \ref{eq:U-Neumann}) for the 
%case in which the horizontal plane $z''_{0}$ crosses the true 3D sources. 
%In this case, however, the potential field $U$ 
%(equation \ref{eq:scalar-potential-U}) cannot be exactly
%reproduced, but approximated to a predefined error level $\zeta$.
%
%PAREI AQUI 2